\scanpage{42}
\mspage{-VI-6-}
\hfill n°383

indices $1,\ldots,n$ ; soit $\pi$ la permutation de $\{1,2,\ldots,n\}$ correspondante ; posons $s_i=s_{H_i}$. Notre assertion est vraie dans les cas suivants :

a) $\pi$ est la permutation circulaire $(1,2,\ldots,n)$ : en effet on a
\[
s_2s_3\ldots s_ns_1=s_1^{-1}(s_1s_2\ldots s_n)s_1\ ;
\]

b) $\pi$ est la transposition $(i,i+1)$ de deux indices $i,i+1$ (ici on identifie $1$ et $n+1$) tels que $s_is_{i+1}=s_{i+1}s_i$ ; c'est, en effet, trivial. Notons que la relation $s_is_{i+1}=s_{i+1}s_i$ revient à dire que $H_i$ et $H_{i+1}$ sont perpendiculaires, c'est-à-dire qu'on a $(e_i|e_{i+1})=0$. Donc, d'après le lemme 5 et la prop.1, il existe un couple d'indices $(i,i+1)$ tels que $s_is_{i+1}=s_{i+1}s_i$.

Il reste donc à montrer que le sous-groupe $\underline{T}$ du groupe symétrique $\underline{S}_n$, engendré par la permutation circulaire $(1,\ldots,n)$ et la transposition $(i,i+1)$ est égal à $\underline{S}_n$. Or, comme $(j,j+1)=(1,\ldots,n)(j-1,j)(1,\ldots,n)^{-1}$ (pour $j=1,\ldots,n$ ; on identifie aussi $0$ et $n$), on voit par récurrence que toute transposition de la forme $(j,j+1)$ appartient à $\underline{T}$. D'autre part, pour tout $\pi\in\underline{S}_n$, la permutation $\pi'=(1,2)(2,3)\ldots(\pi(1)-1,\pi(1))\pi$ laisse l'élément $1$ invariant ; ceci démontre, par récurrence sur $n$, que $\underline{S}_n$ est engendré par les transpositions $(j,j+1)$. On a donc bien $\underline{S}_n=\underline{T}$. C.Q.F.D.

Il résulte de la prop.2 que toutes les transformations de \unsure{poulcrier} de $G$ ont le même ordre $h=h(G)$. Ce nombre s'appelle le \uline{nombre de Coxeter} de $G$.
\mspage{41}
